\chapter{Conclusion and Future Scope}

\section{Conclusion}

This project has successfully demonstrated the design and implementation of a secure, practical digital signature system that addresses real-world challenges of document forgery and unauthorized modification in digital environments. By combining RSA-2048 cryptographic signing with SHA-256 integrity hashing within an accessible web-based interface, the Document Signer System bridges the gap between the theoretical security guarantees of public-key cryptography and the practical needs of users who require straightforward document authentication tools.

The system accomplishes its primary objectives of implementing secure digital signatures, ensuring data integrity through tamper detection, providing an intuitive user interface, and maintaining a well-structured database for all system entities. The complete workflow from user registration and key generation through document upload, signing, and verification has been implemented, tested, and documented with supporting diagrams, screenshots, and test results.

Through this implementation, the development team has gained substantial practical experience in software engineering, applied cryptography, database design, and system architecture. The project validates the feasibility of building enterprise-grade security features using open-source tools and frameworks, demonstrating that robust document authentication does not require expensive commercial platforms or specialized hardware. The transparent cryptographic pipeline, where all operations are implemented in a dedicated utility module, provides educational value and auditability that proprietary solutions typically do not offer.

The testing results presented in Chapter~7 confirm that the system operates reliably across all primary use cases and correctly identifies both valid signatures and tampered documents. The academic report, system diagrams, and application screenshots together provide a complete project submission suitable for final-year evaluation under MAKAUT CS781 Project II requirements.

\section{Limitations}

The current implementation stores private keys in the database without additional encryption, which is acceptable for academic demonstration but would require enhanced key protection in production deployment. The system supports single-signer workflows only and does not implement certificate authority integration or legally certified digital signatures. SQLite is used as the database engine for simplicity, though production environments would benefit from MySQL or PostgreSQL with proper backup and access control mechanisms. These limitations are acknowledged and addressed as future enhancements rather than deficiencies in the core project scope.

\section{Future Scope}

Several enhancements have been identified that would extend the system's capabilities for production deployment. Web interface enhancement would add features such as drag-and-drop file upload, inline PDF document preview, signature visualization overlays, and responsive mobile layouts optimized for smartphones and tablets. Multi-document batch signing would enable users to sign multiple documents in a single operation, improving efficiency for organizations processing large document volumes.

Cloud deployment capability would involve containerizing the application with Docker, configuring MySQL or PostgreSQL for production database requirements, and deploying to cloud platforms such as AWS, Azure, or Google Cloud with proper SSL/TLS certificate management. Mobile application development would extend the signing and verification capabilities to Android and iOS platforms, enabling document authentication from mobile devices in field environments.

Integration with established public key infrastructure and certificate authorities would provide legally recognized digital signatures compliant with regulations such as the Information Technology Act 2000 and the eIDAS framework. Additional cryptographic algorithm support including ECDSA would offer alternative signing schemes with smaller key sizes and faster computation. Blockchain-based signature anchoring could provide immutable timestamping of signed documents for enhanced non-repudiation guarantees.

\section{Final Remarks}

The Document Signer System represents a meaningful contribution to the practical understanding of digital signature technology in an academic setting. It demonstrates that secure document authentication can be implemented using modern web development practices and established cryptographic standards, while remaining accessible to users without specialized security knowledge. The project provides a solid foundation upon which future enhancements can be built, and it serves as an effective demonstration of the integration of cryptography, software engineering, and user-centered design in a single cohesive application.
